% Part: first-order-logic
% Chapter: introduction
% Section: first-order-logic

\documentclass[../../../include/open-logic-section]{subfiles}

\begin{document}

\olfileid{fol}{int}{fol}

\olsection{منطق الرتبة الأولى}

الغالب أن القارئ قد عرف منطق الرتبة الأولى في دراسته الأولى للمنطق
الصوري.\footnote{وعلى هذه المعرفة نبني عرضنا إلى حد بعيد!{} فمن
  احتاج إليها رجع إلى كتاب تمهيدي، مثل \emph{forall x}
  \citep{Magnus2021}.} وقد يكون عرفه باسم «منطق الكم» أو «منطق
المحمولات». وهو في أصله لغة صورية: لها معجم معلوم، وتؤلف عباراتها
من سلاسل عناصره. ولا تعد كل سلسلة عبارة مقبولة؛ بل للمقبول منها
أقسام: الحدود، وال!!{formula}s، وال!!{sentence}s. والذي يعنينا
خاصة هو ال!!{sentence}s، فإنها تقوم في منطق الرتبة الأولى مقام
الجمل في اللغة الطبيعية. ومن ثم تعرض فيها مسائل المنطقي، كأن نسأل:
\begin{itemize}
    \item هل $!B$ لازمة منطقيًا عن~$!A$؟
    \item أَ$!A$ صادقة منطقيًا، أم كاذبة منطقيًا، أم عارضة؟
    \item أَ$!A$ و$!B$ متكافئتان؟
\end{itemize}

مرجع هذه المسائل أولًا إلى «معنى» ال!!{sentence}s. فالفيلسوف إذا
سأل: أتلزم $!B$ منطقيًا عن~$!A$؟ قصد إلى التمييز بين أمرين:
أيمكن أن تصدق $!A$ في حالة تكذب فيها~$!B$، فلا تكون $!B$ لازمة
عن~$!A$؟ أم كل حالة تصدق فيها $!A$ تصدق فيها~$!B$، فتكون $!B$
لازمة عن~$!A$؟ ولم نبيّن بعد ما «الحالة»؛ وبيانها من شأن
\emph{الدلالة}. فدلالة منطق الرتبة الأولى تضبط ضبطًا رياضيًا فكرة
«الحالة» الحدسية عند الفيلسوف، وتضبط معها معنى كون
!!a{sentence}~$!A$ \emph{صادقة في} تلك الحالة. وما نضعه لضبط
الحالة نسميه !!a{structure}، والعلاقة التي نضبط بها «الصدق في»
نسميها \emph{الإرضاء}. فسنعرف «إرضاء $!A$ في~$\Struct{M}$»،
ورمزه $\Sat{\OLId{latin-looped}{M}}{!A}$، حيث ال!!{sentence}s~$!A$ وال!!{structure}s~$\Struct{M}$.
وبعد ذلك نحد سائر الألفاظ الدلالية، مثل «تلزم عن» و«صادقة
منطقيًا»، حدودًا دقيقة. تتيح لنا هذه الحدود أن نفصل رياضيًا في
قولنا، مثلًا:
$\lforall[\OLId{latin-ordinary}{x}][(!A(\OLId{latin-ordinary}{x}) \lif !B(\OLId{latin-ordinary}{x}))],
\lexists[\OLId{latin-ordinary}{x}][!A(\OLId{latin-ordinary}{x})] \Entails \lexists[\OLId{latin-ordinary}{x}][!B(\OLId{latin-ordinary}{x})]$.
والجواب نعم. ومن تمرس بترميز اللغة الطبيعية عرف أن هذا ترميز لحجة
مثل: «كل النمل حشرات، ويوجد نمل، فإذن توجد حشرات». وصلاحية هذه
الحجة ظاهرة؛ فلا بد أن يوافقها تعريفنا الرياضي للزوم في اللغة
الصورية.

ولعل القارئ عرف أيضًا من دراسته الأولى \emph{!!{derivation}s}.
فالمقرر التمهيدي في المنطق الصوري يدرب على إيجاد هذه
ال!!{derivation}s، وقد يكون من أمثلته !!a{derivation} للاستلزام
المتقدم. ولعرض ال!!{derivation}s طرائق كثيرة، منها «الاستنباط
الطبيعي» و«أشجار الصدق» وغيرهما. وإنما توضع أنساق ال!!{derivation}
لتكون أدوات للفصل في الأسئلة السابقة. فمثلًا، !!a{derivation}
بالاستنباط الطبيعي مقدمته
$\lforall[\OLId{latin-ordinary}{x}][(!A(\OLId{latin-ordinary}{x}) \lif !B(\OLId{latin-ordinary}{x}))]$ و$\lexists[\OLId{latin-ordinary}{x}][!A(\OLId{latin-ordinary}{x})]$،
ونتيجته، أي سطره الأخير، $\lexists[\OLId{latin-ordinary}{x}][!B(\OLId{latin-ordinary}{x})]$، \emph{يثبت} أن
$\lexists[\OLId{latin-ordinary}{x}][!B(\OLId{latin-ordinary}{x})]$ لازمة منطقيًا عن
$\lforall[\OLId{latin-ordinary}{x}][(!A(\OLId{latin-ordinary}{x}) \lif !B(\OLId{latin-ordinary}{x}))]$ و$\lexists[\OLId{latin-ordinary}{x}][!A(\OLId{latin-ordinary}{x})]$.

فما وجه هذا الإثبات؟ أنساق ال!!{derivation} في ظاهرها لا تتعرض
للدلالة؛ فإن إقامة !!a{derivation} صوري إنما هي ترتيب رموز بمقتضى
قواعد، ولا ذكر فيها للحالات ولا للصدق فيها. فلا تثبت الصلة بين
أنساق ال!!{derivation} والدلالة إلا ببحث فوق منطقي. نحتاج، مثلًا،
إلى برهان رياضي يثبت أن إمكان !!a{derivation} صوري
لـ~$\lexists[\OLId{latin-ordinary}{x}][!B(\OLId{latin-ordinary}{x})]$ من
$\lforall[\OLId{latin-ordinary}{x}][(!A(\OLId{latin-ordinary}{x}) \lif !B(\OLId{latin-ordinary}{x}))]$ و$\lexists[\OLId{latin-ordinary}{x}][!A(\OLId{latin-ordinary}{x})]$
يكافئ كون $\lforall[\OLId{latin-ordinary}{x}][(!A(\OLId{latin-ordinary}{x}) \lif !B(\OLId{latin-ordinary}{x}))]$ و
$\lexists[\OLId{latin-ordinary}{x}][!A(\OLId{latin-ordinary}{x})]$ مستلزمتين معًا~$\lexists[\OLId{latin-ordinary}{x}][!B(\OLId{latin-ordinary}{x})]$.
ولا نشرع في ذلك حتى نفرغ من العمل الدقيق اللازم لضبط النحو والدلالة.

\end{document}
